\documentclass[10pt,a4paper]{amsart}
\usepackage[margin=19mm]{geometry}
\usepackage{amsmath,amssymb,mathtools}
\usepackage[scr=rsfs,cal=euler]{mathalfa}
\usepackage{xfrac}
\usepackage[unicode,hidelinks,bookmarks=false]{hyperref}
\newcommand{\ie}{i.e.\ }
\newcommand{\cf}{cf.\ }
\newcommand{\resp}{respectively\ }
\newcommand{\Subsection}{\subsection}
\providecommand{\nobreakdash}{\nobreak\hbox{-}}
\setlength{\emergencystretch}{2em}
\multlinegap0pt
\usepackage{amssymb}
\usepackage{mathtools}
\usepackage{dsfont}
\usepackage{tikz}
\newtheorem{lem}{Lemma}
\newtheorem{prop}{Proposition}
\newtheorem{thm}{Theorem}
\def\A{\mathcal A} 
\def\C{\mathcal C}
\newcommand{\End}{\operatorname{End}}
\newcommand{\Hom}{\operatorname{Hom}}
\newcommand{\Inj}{\operatorname{Inj}}
\def\La{\Lambda}
\newcommand{\Rad}{\operatorname{Rad}}
\newcommand{\Spec}{\operatorname{Spec}}
\def\U{\mathcal U}
\def\bfB{\mathbf B} 
\def\bfD{\mathbf D} 
\def\bfL{\mathbf L} 
\def\bfS{\mathbf S}
\newcommand{\iso}{\xrightarrow{\raisebox{-.4ex}[0ex][0ex]{$\scriptstyle{\sim}$}}}
\newcommand{\longiso}{\xrightarrow{\ \raisebox{-.4ex}[0ex][0ex]{$\scriptstyle{\sim}$}\ }}
\newcommand{\supp}{\operatorname{supp}}
\title{The Boolean spectrum of a Grothendieck category}
\author{Henning Krause}
\date{Source: arXiv:2411.16309v3 (CC BY 4.0)}
\begin{document}
\maketitle

\noindent\emph{Source and licence.} The Boolean spectrum of a Grothendieck category. Version arXiv:2411.16309v3 (CC BY 4.0), distributed by arXiv under the Creative Commons Attribution 4.0 International licence, http://creativecommons.org/licenses/by/4.0/, as recorded in the arXiv licence metadata for this exact version and shown on the abstract page https://arxiv.org/abs/2411.16309v3. Full licence notice: \texttt{licenses/arxiv-2411.16309-CC-BY-4.0.txt}.\par\smallskip

\noindent\textit{Editorial scope.} Counted unit: the article's thm th:decomp (source0.tex lines 1495--1500). The article's complete original proof of it is delivered with it (source0.tex lines 1501--1529, one \texttt{proof} environment of 29 lines). No mathematics is written, repaired or re-derived: the statement and the proof are the article's own lines, in the article's own notation, and no retained line is substituted or re-flowed.  Retained uncounted as the source-local context the proof needs: lines 547--549; lines 566--570; lines 738--741; lines 822--829; lines 1093--1098.  Nothing was omitted from any retained span, so the package carries no omission record.  No retained line contains a citation, so the wrapper carries no bibliography and no bibliography entry.  No retained line contains a figure environment, an image-inclusion command or drawing code, so no image is delivered and the package declares no asset dependency.  Editorial formatting only, in addition: an AMS \texttt{amsart} preamble that restates the article's own declarations and macros used by the retained lines, namely the environments \texttt{lem}, \texttt{prop}, \texttt{thm} on separate counters, together with the standard packages those lines need.  The retained environments print in this document as Lemma 1, Proposition 1, Theorem 1 in order of appearance, whereas the article prints the same items with the numbering of its own full document.  The complete native source of the article as distributed in the arXiv e-print for this exact version is bundled unchanged as \texttt{sources/169055/source0.tex}.
\begin{equation}\label{eq:Spec}
  (\Inj\A)/\Rad(\Inj\A)\longiso\Spec\A,
\end{equation}

\begin{lem}\label{le:spec-decomp-inj}
  Let $X\in\A$ be an injective object and $P(X)=Y_1\oplus Y_2$ a
  decomposition. Then there exists a decomposition $X=X_1\oplus X_2$
  such that $Y_i=P(X_i)$ for $i=1,2$.
\end{lem}

\begin{prop}\label{pr:Boole}
  Let $\A$ be a spectral Grothendieck category.  Then the lattice
  $\bfL(\A)$ is a complete Boolean lattice.
\end{prop}

\begin{lem}\label{le:centre}
  Let $\A$ be a spectral Grothendieck category. Then the assignment
  $\C\mapsto e_\C$ induces a lattice isomorphism
  \[\bfL(\A)\longiso\bfB(Z(\A)).\]
  Moreover, when $G\in\A$ is a generator with $\La=\End(G)$, then
   \[x\longmapsto\{X\in\A\mid \Hom(G,X)x=\Hom(G,X)\}\] induces a lattice isomorphism
  \[\bfB(\La)\longiso\bfL(\A).\]
\end{lem}

\begin{thm}\label{th:spectral}
  Let $\A$ be a Grothendieck category. The assignment
  \[\C\longmapsto\supp(\C)=P(\C)\] induces a lattice isomorphism
\[\{\C\subseteq\A\mid \C\text{ essentially closed}\}\longiso\bfS(\A).\]
The inverse map is given by  $\U\mapsto P^{-1}(\U)$.
\end{thm}

\begin{thm}\label{th:decomp}
  Given an injective object $X$ in $\A$, the partially ordered set $\bfD(X)$ of direct
  summands is a complete Boolean lattice. Moreover, there is a
  canonical lattice isomorphism
\[\bfD(X)\longiso\bfB(\End(X)/J(\End(X))).\]
\end{thm}

\begin{proof}
We claim that the assignment $U\mapsto \langle P(U)\rangle= P(\langle
U\rangle)$ induces a lattice isomorphism
\begin{equation}\label{eq:decomp}
  \bfD(X)\longiso\{\U\in\bfS(\A)\mid \U\subseteq\supp(X)\}.
\end{equation}  
The inverse map sends a localising subcategory $\U\subseteq\Spec\A$ to
the equivalence class of a direct summand $U\subseteq X$ such that
$P(U)=t_\U(P(X))$; for the existence of $U$ see
Lemma~\ref{le:spec-decomp-inj}. Using Theorem~\ref{th:spectral} it is
easily checked that these assignments are mutually inverse to each
other and preserving the partial order. Thus $\bfD(X)$ is a Boolean
lattice, since $\bfS(\A)$ has this property by
Proposition~\ref{pr:Boole}. Moreover, the above argument shows that
$P\colon\A\to\Spec\A$ induces a lattice  isomorphism
\begin{equation}\label{eq:decomp1}
  \bfD(X)\longiso\bfD(P(X)).
\end{equation}

For the isomorphism with the lattice of central idempotents we
consider the Grothendieck category
\[\A_X:=\langle P(X)\rangle\subseteq\Spec\A\] which is generated by
$P(X)$. The functor $P$ induces an isomorphism
\[\End(X)/J(\End(X))\longiso \End(P(X))\] by \eqref{eq:Spec}. This
 yields the following sequence of isomorphisms
 \[\bfD(X)\iso\bfD(P(X))\iso\bfL(\A_X)\iso\bfB(\End(X)/J(\End(X))),\]
 where the first is \eqref{eq:decomp1}, the second is \eqref{eq:decomp}, and the last comes from
 Lemma~\ref{le:centre}.
\end{proof}

\end{document}
